\section*{Limitations}
\label{sec:limitations}

This work has several limitations that we want to make explicit before they become reviewer surprises.

\paragraph{Statistical power.} Our $T_{\text{eval}}$ evaluation runs at $K{=}3$ seeds, giving $n{=}54$ per Tier-B arm and $n{=}51$ per Tier-C arm. The directional deltas we report---$+9.3$~pp for B2$-$B1, $+3.7$~pp for B3$-$B2---are not significant at $\alpha=0.05$ on their own ($z\approx 0.4$--$1.0$). The signal is consistent across two independently seeded comparisons and concentrates on the largest-$n$ Tier-B app in a methodologically interpretable way, but a wider $K_{\text{eval}}$ sweep would be needed to clear the formal significance bar.

\paragraph{Single backbone.} All numbers are on Qwen3-VL-8B-Instruct. Whether the structural priors and the EvoFSM-RL loop carry over to a stronger or larger vision-language model is unanswered. The choice was made on the joint constraint of fitting LoRA training on a single 80\,GB card and matching the M3A two-phase loop introduced with AndroidWorld; relaxing either constraint changes which backbones become reachable.

\paragraph{LAYER 1 is rediscovered in-context.} The full algorithmic description specifies a target-app cold-start phase that uses the pretrained policy to roll out a few exploration trajectories and synthesize a target-app LAYER 1 before adaptation begins. Our current implementation collapses that phase to a single category look-up and lets the agent rediscover LAYER 1 from screenshots during the B3/B4 loop itself. This works because the visual encoder is strong; it works less well in cases like \texttt{chrome/BrowserMultiply}, where the failure mode is upstream of any FSM (the agent emits a valid but mislocated \texttt{open\_app: ``File Manager''} action repeatedly until step exhaustion). A target-app LAYER 1 extraction phase is a natural follow-up.

\paragraph{Sparse-template apps.} Three Tier-B apps and several Tier-C apps have $|T_{\text{eval}}|\le 3$. Per-app deltas on these apps swing on single-seed flips and should be read as qualitative checks rather than rates. The aggregate Tier-B number is robust to removing the small-$n$ rows individually but fragile to removing the dominant $n{=}18$ app (\texttt{system\_settings}); we report the per-app breakdown alongside every aggregate.

\paragraph{The base-policy precondition.} GRPO's within-group baseline degenerates to zero when every rollout on a $(F_i, t)$ cell fails uniformly. This happens on apps where the base policy's success rate on $T_{\text{adapt}}$ is near zero (\texttt{simple\_calendar\_pro}, \texttt{pro\_expense}, \texttt{retro\_music}, and most of Tier-C), and it leaves the LoRA update with no positive anchor to imitate from. On apps where the base policy succeeds at meaningful rates (\texttt{system\_settings}), the joint update improves further. This precondition is shared with all RL-from-rollouts methods, but it is more visible in our deployment-time TTA setting than in large-budget pretraining; we discuss reward shaping and demonstration warm-start as the natural next steps.

\paragraph{B4 evaluation is the central outstanding piece.} The four-rung ablation needs the B4 row across both tiers to close. The training pipeline is in place---grouping, normalization, optimizer, the Phase 1 / Phase 3 split, and a category-matched warm start are all wired---but the end-to-end B4 success-rate number on $T_{\text{eval}}$ is the missing column in this version of the paper.
